\chapter{Simple I/O}

An important part of any computer program is to handle input and output. In our examples so far, we have already used the two most common Fortran constructs for this: \texttt{read} and \texttt{write}. Fortran I/O can be quite complicated, so we will only describe some simpler cases in this tutorial.

\section{Read and write}

\texttt{Read} is used for input, while \texttt{write} is used for output. A simple form is

\begin{lstlisting}[style=fortran77]
      read (unit no, format no) list-of-variables
      write(unit no, format no) list-of-variables
\end{lstlisting}

The unit number can refer to either standard input, standard output, or a file. The format number refers to a label for a format statement, which will be described in the next lesson.

It is possible to simplify these statements further by using asterisks (*) for some arguments, like we have done in most of our examples so far. This is sometimes called \emph{list directed} read/write.

\begin{lstlisting}[style=fortran77]
      read (*,*) list-of-variables
      write(*,*) list-of-variables
\end{lstlisting}

The first statement will read values from the standard input and assign the values to the variables in the variable list, while the second one writes to the standard output.

\section{Examples}

Here is a code segment from a Fortran program:

\begin{lstlisting}[style=fortran77]
      integer m, n
      real x, y, z(10)

      read(*,*) m, n 
      read(*,*) x, y
      read(*,*) z
\end{lstlisting}

We give the input through standard input (possibly through a data file directed to standard input, as in \texttt{./prog < data.txt}). A data file consists of \emph{records} according to traditional Fortran terminology; for a text file, a record is simply a line. Each \texttt{read} statement starts reading at the beginning of a new record. Within a record, the values are separated by blanks or commas. Hence a legal input to the program above would be:

\begin{lstlisting}[style=fortran77]
   -1  100
  -1.0 1e+2
  1.0 2.0 3.0 4.0 5.0 6.0 7.0 8.0 9.0 10.0
\end{lstlisting}

Or, we could add commas as separators:

\begin{lstlisting}[style=fortran77]
   -1, 100
 -1.0, 1e+2
  1.0, 2.0, 3.0, 4.0, 5.0, 6.0, 7.0, 8.0, 9.0, 10.0
\end{lstlisting}

Note that Fortran 77 input is line sensitive, so it is important not to have extra input elements (fields) on a line (record). For example, if we gave the first four inputs all on one line as

\begin{lstlisting}[style=fortran77]
   -1, 100, -1.0, 1e+2
  1.0, 2.0, 3.0, 4.0, 5.0, 6.0, 7.0, 8.0, 9.0, 10.0
\end{lstlisting}

then m and n would be assigned the values -1 and 100 respectively, but the last two values would be discarded, because the next \texttt{read} starts on a new line. So x and y would be assigned the values 1.0 and 2.0, ignoring the rest of the second line. Then the third \texttt{read} finds no more input for z, and the program stops with an error:

\begin{lstlisting}[style=output]
Fortran runtime error: End of file
\end{lstlisting}

(In Chapter 15 we saw how to handle the end of the input with \texttt{END=} or \texttt{IOSTAT=} instead of letting the program stop.)

If there are too few inputs on a line then the next line will be read. For example

\begin{lstlisting}[style=fortran77]
   -1
  100
 -1.0
 1e+2
  1.0, 2.0, 3.0, 4.0, 5.0, 6.0, 7.0, 8.0, 9.0, 10.0
\end{lstlisting}

would produce the same results as the first two examples.

\section{Other versions}

For simple list-directed I/O it is possible to use the alternate syntax

\begin{lstlisting}[style=fortran77]
      read  *, list-of-variables
      print *, list-of-variables
\end{lstlisting}

which has the same meaning as the list-directed read and write statements described earlier. This version always reads/writes to standard input/output so the * corresponds to the format.

\section{Implied do-loops}

To read or write many elements of an array in a single statement, Fortran has the \emph{implied do-loop}: a list of items, followed by a loop control as in a do statement, all between parentheses:

\begin{lstlisting}[style=fortran77]
      read(*,*) (x(i), i = 1, n)
      write(*,*) (x(i), i = 1, n, 2)
\end{lstlisting}

The first statement reads \texttt{x(1)}, \texttt{x(2)}, \ldots, \texttt{x(n)}; since it is a single read statement, the n values can be distributed over the input lines in any way. (A do loop with a read statement inside would instead start a new line for every element.) The second statement writes the elements with odd index. Implied do-loops can be nested, and are often combined with an ordinary do loop to print a matrix row by row:

\begin{lstlisting}[style=fortran77]
      program implid
      implicit none
      integer n, i, j
      real x(10), A(3,3)

      read(*,*) n
      read(*,*) (x(i), i = 1, n)
      write(*,*) 'x:', (x(i), i = 1, n)
      write(*,*) 'Odd positions:', (x(i), i = 1, n, 2)

      do 20 j = 1, 3
         do 10 i = 1, 3
            A(i,j) = real(10*i + j)
   10    continue
   20 continue
      write(*,*) 'A, row by row:'
      do 30 i = 1, 3
         write(*,*) (A(i,j), j = 1, 3)
   30 continue
      write(*,*) 'A in memory order:'
      write(*,*) ((A(i,j), i = 1, 3), j = 1, 3)
      end
\end{lstlisting}

With the input

\begin{lstlisting}[style=fortran77]
5
1.5 2.5 3.5
4.5 5.5
\end{lstlisting}

the program prints (the last line is shown here broken in two)

\begin{lstlisting}[style=output]
 x:   1.50000000       2.50000000       3.50000000       4.50000000       5.50000000
 Odd positions:   1.50000000       3.50000000       5.50000000
 A, row by row:
   11.0000000       12.0000000       13.0000000
   21.0000000       22.0000000       23.0000000
   31.0000000       32.0000000       33.0000000
 A in memory order:
   11.0000000       21.0000000       31.0000000       12.0000000       22.0000000
   32.0000000       13.0000000       23.0000000       33.0000000
\end{lstlisting}

In the nested implied loop the \emph{inner} loop (over \texttt{i}) runs fastest, so the elements come out in the order in which they are stored in memory, column by column. Writing just the name of the array, \texttt{write(*,*) A}, gives the same result.

\section{Exercises}

Solutions are in the directory \texttt{code/exercises/ch16} that accompanies this tutorial.

\begin{enumerate}
    \item What values do the variables receive if the program

\begin{lstlisting}[style=fortran77]
      integer a, b, c, d, e
      read (*,*) a, b
      read (*,*) c, d
      read (*,*) e
\end{lstlisting}

    reads the input below? Answer first without a computer.

\begin{lstlisting}[style=fortran77]
1 2 3
4
5
6 7
\end{lstlisting}

    \item Write a program that reads an integer $n \le 100$ and then $n$ integers, distributed over the input lines in any way, with an implied do-loop. The program should print their sum, and the numbers in reverse order (use an implied do-loop with a negative increment).
    \item Write a program that reads a $3 \times 3$ matrix row by row, one row on each line, and prints its transpose row by row.
    \item Run the program \texttt{checks/eof\_extra.f} that accompanies this tutorial with the input \texttt{checks/eof\_extra.in} (the case with the first four values on one line). Explain the error message. Which values would \texttt{x} and \texttt{y} receive?
\end{enumerate}
